The angle sum is 1260, so the arithmetic progression is $140-4d,\ldots,140+4d$. Integer angles, distinctness and convexity give d=1 through nine. These are nine angle progressions, but fixed angles do not determine a polygon up to similarity: side lengths can vary continuously while preserving the angles and convexity. Thus there are $\boxed{\text{infinitely many}}$ non-similar nonagons.
\begin{asy}
import graph;import geometry;usepackage("amsmath");size(7cm);
pair[] V={(0,0),(1.0470723783,0),(1.83229369268,0.732230721622),(1.92480106891,1.7895948703),(1.32592684303,2.61387452744),(0.386446675922,2.84811324121),(-0.459164958772,2.45379806013),(-0.880186425021,1.62749690743),(-0.714955231069,0.690424241133)};path p=V[0];for(int i=1;i<9;++i)p=p--V[i];draw(p--cycle);for(int i=0;i<9;++i)label("$"+string(136+i)+"^\circ$",V[i],unit(V[i]-(0,1.5)));
\end{asy}

\begin{asy}
import graph;import geometry;usepackage("amsmath");size(7cm);
pair[] V={(0,0),(1.433552912,0),(2.15795862634,0.675519257106),(2.24958689531,1.72283516386),(1.61158466717,2.60096989589),(0.565274315585,2.86184436621),(-0.373655615449,2.42401414889),(-0.818152548096,1.55163979922),(-0.656311296592,0.633792451898)};path p=V[0];for(int i=1;i<9;++i)p=p--V[i];draw(p--cycle);for(int i=0;i<9;++i)label("$"+string(136+i)+"^\circ$",V[i],unit(V[i]-(0,1.5)));
\end{asy}

